\chapter{战斗有效属性与状态变量}

若不先把战斗里真正参与判定的变量层次理清，后文关于乱敏、命中、会心、法术命中与属性状态的公式就很容易混成同一层。对当前 \code{gmsv} 而言，最重要的区别不是“面板属性和战斗属性谁更大”，而是“同一个中文词在源码里究竟落在哪个工作变量上”。本章因此先把人物、宠物、敌人在战斗中的有效属性层次、身份差异和常见状态变量整理出来。

\section{原始属性、固定属性与战斗属性}

从源码结构看，战斗中至少存在三层不同的属性口径。第一层是\emph{原始属性}，例如人物的 \code{CHAR_STR}、\code{CHAR_TOUGH}、\code{CHAR_DEX}、\code{CHAR_LUCK}，以及宠物由成长与等级生成后的显示攻击、防御、敏捷。它们更接近“回合开始前已经确定的基础值”。第二层是\emph{固定属性}，例如 \code{CHAR_WORKFIXSTR}、\code{CHAR_WORKFIXTOUGH}、\code{CHAR_WORKFIXDEX}、\code{CHAR_WORKFIXLUCK}，它们是在装备、套装、部分状态和战斗初始化结算之后写入的当前回合工作值。第三层则是\emph{战斗属性}，例如 \code{CHAR_WORKATTACKPOWER}、\code{CHAR_WORKDEFENCEPOWER}、\code{CHAR_WORKQUICK}，它们是在固定属性之上，再经过本回合的攻防速修正后得到的最终判定量。

把这一过程抽象写成函数链，可以记为
\begin{equation}
    X_{\mathrm{raw}}
    \xrightarrow{\;\mathcal{F}\;}
    X_{\mathrm{fix}}
    \xrightarrow{\;\mathcal{T}\;}
    X_{\mathrm{battle}},
\end{equation}
其中 $\mathcal{F}$ 表示装备、状态和战斗初始化引起的固定化过程，$\mathcal{T}$ 表示 \code{BATTLE_TurnParam()} 这一类“按本回合修正生成攻、防、速工作值”的过程。以敏捷为例，更具体的关系可以理解为
\begin{equation}
    \dex_{\mathrm{raw}} \longrightarrow \dex_{\mathrm{fix}}=\code{WORKFIXDEX}
    \longrightarrow Q=\code{WORKQUICK}.
\end{equation}
因此，玩家常说的“敏捷”在不同公式里并不总是指同一个量：行动顺序主要看 $Q$，而物理回避、会心、反击更多直接看 $\dex_{\mathrm{fix}}$。

幸运也存在类似分层，但其口径更容易混淆。物理回避、会心、反击常读的是固定幸运 \code{CHAR_WORKFIXLUCK}；而普通属性法术与职业法术命中，核心却仍直接读取原始幸运 \code{CHAR_LUCK}。这意味着“同样叫幸运”，在物理系统和法术系统里进入公式的位置并不一致。对玩家研究而言，这种变量层次差异比“某个系数是 $0.8$ 还是 $0.6$”更基础。

\section{人物、宠物与敌人的身份差异}

战斗公式中很多修正并不是由技能决定，而是先由单位身份决定。当前实现至少区分人物、宠物和敌方单位三类身份，它们会通过不同缩放系数进入同一条概率公式。以物理回避与会心为例，源码会根据“谁打谁”的组合，对攻方或守方敏捷先乘一次 $0.8$ 或 $0.6$ 再截断。其含义不是“宠物永远比较慢”或“人物永远更容易命中”，而是不同身份组合使用了不同的标准化口径。

从当前已核对的公式可抽象出两类最常见的身份差异。第一类是\emph{敏捷缩放差异}：人物打非人物、非人物打人物、敌人与宠物交战时，攻守双方敏捷会先做不同折扣，再进入回避或会心的平方根结构。第二类是\emph{替代变量差异}：当法术目标不是玩家时，系统往往不再读取幸运与抗魔装备，而改用目标等级作为近似项；当目标是宠物时，幸运在主 PVP 口径下通常也可近似视作 $0$。因此，“某技能打宠和打人感觉不一样”并不总来自技能本身，很多时候只是因为输入变量已经换了一套。

这种身份差异对后续章节有两个直接影响。其一，命中、回避、会心等热力图若只画“人物对人物”，不能直接外推到“人物对宠”或“宠对敌”；其二，宠物养成的价值不能只用人物 PVP 公式衡量，因为宠物本身落到战斗里时，输入层级和身份缩放都可能发生变化。也正因为此，后文在讨论“从宠物成长到战斗能力”时，会把人物、宠物、骑宠三种口径分开处理。

\section{常见战斗状态变量}

源码中的战斗状态并不都以“一个技能对应一个布尔量”的方式存在。更准确地说，它们是一组分散在 \code{CHAR_WORK*}、\code{CHAR_WORKBATTLEFLG}、\code{CHAR_WORKBATTLECOM*} 和战场全局结构中的状态槽。下面先列出对后文最关键的一组变量。

首先是若干会直接改变判定入口的布尔标记。防御动作 \code{BATTLE_COM_GUARD} 会让普通回避在最前面直接失败；\code{CHAR_BATTLEFLG_NODUCK} 表示“不能闪避”；\code{CHAR_BATTLEFLG_REVERSE} 则表示单位当前处于属性反转状态。其次是若干会改变具体参数的工作变量，例如 \code{CHAR_WORKBATTLECOM3} 会被当前动作复用，高 16 位和低 16 位在不同技能中可分别表示技能等级、子类型、阵列或额外修正值；\code{CHAR_MYSKILLHIT} 与 \code{CHAR_MYSKILLHIT_NUM} 则对应某些“提高命中”类主动状态。

这里还需要记住一个与宠物相关的战斗标记：\code{CHAR_BATTLEFLG_AIBAD}。本文把它称为宠物“\textbf{nono / 忠诚失控}”标记，它在 \code{BATTLE_PetLoyalCheck()} 中按回合清空后重新判断，表示该宠物这一回合没有按玩家原始指令行动；它会进一步影响反击、同队误伤保护等后续分支。

在回避系统里，还需要特别区分两套名字很像、但进入位置完全不同的状态。其一是 \code{_PETSKILL_SETDUCK} 相关的宠技回避状态，它主要借助 \code{CHAR_MYSKILLDUCK} 与 \code{CHAR_MYSKILLDUCKPOWER} 两个工作槽，分别记录剩余回合与独立前判概率；只要这层前判成功，整次物理攻击就直接判定为闪避。其二是职业技能 \textbf{回避} 对应的 \code{CHAR_WORK_P_DUCK} 与 \code{CHAR_WORKMOD_P_DUCK}，前者是“本回合是否启用该倍率”的标记，后者是倍率值；它们并不单独生成一次新的闪避事件，而是在普通回避值已经算出后再做乘法放大。

再往下是与属性系统直接相关的一组状态槽。四属性结界分别写在
\code{CHAR_WORKFIXEARTHAT_BOUNDARY}、
\code{CHAR_WORKFIXWATERAT_BOUNDARY}、
\code{CHAR_WORKFIXFIREAT_BOUNDARY}、
\code{CHAR_WORKFIXWINDAT_BOUNDARY}；
属性抗性会写入 \code{CHAR_WORK_F_RESIST}、\code{CHAR_WORK_I_RESIST}、\code{CHAR_WORK_T_RESIST}；场地属性、场地强度与回合数则存放在战场级结构 \code{BattleArray[battleindex]} 中。也就是说，属性系统并不是单一状态，而是一组作用层级不同、刷新逻辑也不同的子状态集合。

\section{状态变量的存储与回合更新}

从实现方式看，当前战斗状态大体可以分成三种。第一种是\emph{单纯标志位}，例如 \code{NODUCK}、\code{REVERSE} 等，通常通过按位与、按位或、按位异或来开关。第二种是\emph{整数槽位}，例如 \code{WORKFIXDEX}、\code{WORKQUICK}、各类抗性、各类命中附加，它们在每个回合的预处理阶段重新写入。第三种是\emph{高低位打包状态}，典型代表就是四属性结界使用的
\begin{equation}
    \code{MAKE2VALUE(power,turn)},
\end{equation}
它用高 16 位存强度、低 16 位存剩余回合。

状态的时间更新方式也并不统一。场地是战场级状态，每回合把 \code{att_count} 减一，减到非正时清空场地属性；四属性结界则是单位级状态，每回合把低 16 位回合数减一，减到 $-1$ 时清空该结界；属性反转则不是“每回合自动衰减”，而是靠技能再次命中时用异或切换开关。

从研究视角看，这些状态变量最重要的共性在于：它们不是平行地叠在最终伤害上，而是分散进入“预处理”“概率判定”“属性修正”“回合衰减”和“死亡收尾”等不同阶段。正因为状态进入位置不同，所谓“同样是减伤”或“同样是降敏”在源码里可能完全不是同一种机制。后续关于乱敏、法术命中、属性反转和骑宠伤害的章节，都会反复用到本章整理的变量层次与状态分类。

